\documentclass[11pt,a4paper]{article}
\usepackage[T1]{fontenc}
\usepackage{lmodern,amsmath,amssymb,amsthm,mathtools,booktabs,microtype}
\usepackage[margin=27mm]{geometry}
\usepackage[colorlinks=true,linkcolor=blue,citecolor=blue,urlcolor=blue]{hyperref}
\setlength{\emergencystretch}{2em}
\newtheorem{theorem}{Theorem}[section]
\newtheorem{lemma}[theorem]{Lemma}
\newtheorem{corollary}[theorem]{Corollary}
\newtheorem{proposition}[theorem]{Proposition}
\newcommand{\Q}{\mathbb Q}\newcommand{\Z}{\mathbb Z}
\newcommand{\A}{\mathbb A}\newcommand{\calV}{\mathcal V}
\DeclareMathOperator{\Tr}{Tr}\DeclareMathOperator{\Res}{Res}
\title{Quadratic obstructions to a compressed elliptic pairing}
\author{}\date{16 November 2026}
\begin{document}\maketitle
\begin{abstract}
A pairing of two elliptic points described by ten field coordinates
admits a spurious one-parameter family at opposite-point tangencies.
Its proposed difference abscissa can vanish or be four times the true
value, including when both coupling coordinates are nonzero.
At an explicit fixed reference, the space of quadratic identities of
the geometric secant image has dimension nineteen, and every identity
vanishes on the spurious family. Exact nonzero minors, together with
monic normal forms, certify this statement in characteristic zero.
Degree-bounded zero closures commute with products, so the obstruction
also applies to arbitrary rational quadrics in the twenty rational
components, including products of conjugate components. Adjoining the
physical rational trace still leaves the false height unrestricted:
the resulting quadratic identity space has dimension fifty-nine.
Different embeddings, genuinely new helper coordinates and restricted
arithmetic witness families are not covered.
\end{abstract}

\section{The compressed pairing}
Let
\begin{equation}\label{eq:constants}
\begin{gathered}
 K=\Q(\rho),\qquad \rho^2=2,\qquad
 e=-1+\frac74\rho,\qquad
 a=e^2=\frac{57}{8}-\frac72\rho,\\
 c=-\frac{87}{41}+\frac{129}{164}\rho,\qquad
 D=(3-\rho)^2=a(a-c^2),\qquad E:y^2=x^3-Dx.
\end{gathered}
\end{equation}
The points
\[
 R_0=(1+\rho,2\rho),\qquad P=2R_0=(a,ac)
\]
are verified by the duplication formula. No assertion about the
Mordell--Weil group is needed below.

A reference in the unit chart of~\cite{Reference} has constants
$u,w,v,r,q,S,T$ satisfying
\begin{equation}\label{eq:reference}
\begin{gathered}
 auw=1,\qquad v\bigl(v-a(u+w)-2c\bigr)=D/a,\\
 r=u/v,\qquad q=w/v,\\
 S=v-c-au,\qquad T=aw-v+c,\\
 rS^2=1-Dr^2,\qquad qT^2=1-Dq^2.
\end{gathered}
\end{equation}
Here $r=1/x(H)$, $q=1/x(H-P)$ and $S=y(H)/x(H)$,
$T=y(H-P)/x(H-P)$. We use references with $r,q,S,T$ nonzero
for the quadratic-identity theorem. The first counterexample only
requires $r,q$ nonzero.

Work at the nonzero input specialization $\beta=1$. Retain the ten
field coordinates
\[
                     (F,X,L,I,G,Y,N,J,A_1,A_2).
\]
The numerator equations and two proposed coupling equations are
\begin{align}
 I&=(L-1)^2,& FX+DF^2&=r^2,&(I+2L-1)X&=1-Dr^2,
                                                        \label{eq:first}\\
 J&=(N-1)^2,& GY+DG^2&=q^2,&(J+2N-1)Y&=1-Dq^2,
                                                        \label{eq:second}\\
 A_1(X+DF)&=A_2(Y+DG),&&&\label{eq:coupling1}\\
 e(wA_2-uA_1)&=SN+TL.&&&\label{eq:coupling2}
\end{align}
The first seven equations are quadratic; \eqref{eq:coupling2} is
linear when the reference is fixed.

\begin{lemma}\label{lem:recovery}
Every solution of \eqref{eq:first} with $L\ne0$ recovers a finite
point $U\in E$ with nonzero abscissa through
\[
 f=F/r=1/x(U),\qquad s=S/L=y(U)/x(U).
\]
The analogous formulas $g=G/q$, $t=T/N$ recover $V$ from
\eqref{eq:second} when $N\ne0$. In particular $F,G\ne0$.
\end{lemma}
\begin{proof}
The equation $FX+DF^2=r^2$ forces $F\ne0$. Since
$I+2L-1=L^2$ and $rS^2=1-Dr^2$, the last equation of
\eqref{eq:first} gives $X=rs^2$. Dividing the middle equation
by $r^2$ gives $fs^2+Df^2=1$. Thus $(1/f,s/f)$ lies on $E$.
The second block is identical.
\end{proof}

On the open set $fgst(g-f)\ne0$, the true secant pairing is the map
\begin{equation}\label{eq:map}
\begin{gathered}
 s^2=f^{-1}-Df,\qquad t^2=g^{-1}-Dg,\\
 F=rf,\quad X=rs^2,\quad L=S/s,\quad I=(L-1)^2,\\
 G=qg,\quad Y=qt^2,\quad N=T/t,\quad J=(N-1)^2,\\
 A_1=ewLN\frac{f(s+t)}{g-f},\qquad
 A_2=euLN\frac{g(s+t)}{g-f}.
\end{gathered}
\end{equation}
Let $\calV\subset\A_K^{10}$ be the Zariski closure of its geometric
image. This definition requires identities on the complete geometric
secant image, not only on a selected arithmetic witness set.

\section{The tangent family and the false height}
\begin{lemma}\label{lem:line}
On the locus recovered in Lemma~\ref{lem:recovery}, put
\[
 x=x(U),\quad y=x(V),\qquad
 A=\frac{euA_1}{LN},\quad B=\frac{ewA_2}{LN}.
\]
Equations \eqref{eq:coupling1}--\eqref{eq:coupling2} are equivalent to
\begin{equation}\label{eq:line}
                         Ax=By,\qquad B-A=s+t.
\end{equation}
If $x\ne y$, they give $A=M-s$, $B=M+t$, where $M$ is the slope
through $U$ and $-V$, and
\begin{equation}\label{eq:height}
                     \frac{A_1A_2}{L^2N^2}=AB=x(U-V).
\end{equation}
If $x=y$, they instead permit every $A=B$ when $V=-U$.
\end{lemma}
\begin{proof}
The point equations give $X+DF=rx$ and $Y+DG=qy$.
Because $wr=uq$ and $e^2uw=1$, the two coupling equations become
\eqref{eq:line}. For $x\ne y$,
\[
 M=\frac{xs+yt}{x-y},\qquad
 A=\frac{y(s+t)}{x-y},\qquad
 B=\frac{x(s+t)}{x-y}.
\]
The curve equations imply
\begin{equation}\label{eq:divided}
                  M(s-t)=x+y-st.
\end{equation}
Consequently $AB=M^2-x-y=x(U-V)$. If $x=y\ne0$,
\eqref{eq:line} says only $A=B$ and $s+t=0$.
\end{proof}
At a non-torsion tangent $V=-U$, the missing divided equation
\eqref{eq:divided} fixes
\[
                         A=B=\frac{x+D/x}{2s}.
\]
The coupling equations alone leave that value free.

\begin{proposition}\label{prop:zero}
At $H=U=R_0$, $V=-R_0$, the coordinate tuple
\begin{equation}\label{eq:zerofixture}
 F=G=3-2\rho,\quad X=Y=40\rho-56,\quad
 L=N=1,\quad I=J=A_1=A_2=0
\end{equation}
satisfies \eqref{eq:first}--\eqref{eq:coupling2}.
Its proposed height in \eqref{eq:height} is zero, while
\[
                         \Tr_{K/\Q}x(U-V)=57/4.
\]
\end{proposition}
\begin{proof}
Here $r=q=\rho-1$, $S=4-2\rho$, $T=-S$.
Set $v=(S+c)/(1-ar)$ and $u=w=vr$; these satisfy
\eqref{eq:reference}. Direct substitution gives every equation.
The actual difference is $2R_0=P$, whose abscissa is $a$.
\end{proof}

For the main obstruction fix $H=2P$. Its reference constants are
\begin{equation}\label{eq:fixedref}
\begin{gathered}
 v=\frac{c^2-2a}{2c},\qquad r=v^{-2},\quad q=a^{-1},\quad
 u=v^{-1},\quad w=v/a,\\
 S=v-c-a/v,\qquad T=c.
\end{gathered}
\end{equation}
They satisfy \eqref{eq:reference}, with all displayed denominators
and $r,q,S,T$ nonzero. Duplication gives $x(2P)=v^2$.

\begin{proposition}\label{prop:family}
For every $\zeta\in K$, the tuple
\begin{equation}\label{eq:family}
\begin{gathered}
 F=r/a,\quad X=rc^2,\quad L=S/c,\quad I=(S/c-1)^2,\\
 G=q/a,\quad Y=qc^2,\quad N=-1,\quad J=4,\\
 A_1=ewLN\zeta,\qquad A_2=euLN\zeta
\end{gathered}
\end{equation}
satisfies all eight equations. Its recovered points are $U=P$,
$V=-P$, and its proposed difference abscissa is $\zeta^2$.
The choice $\zeta=2v$ has $A_1A_2\ne0$ and gives four times
the actual difference abscissa.
\end{proposition}
\begin{proof}
Recovery gives $f=g=1/a$, $s=c$, $t=-c$.
Lemma~\ref{lem:line} permits $A=B=\zeta$, proving the equations.
The normalization $e^2uw=1$ gives
$A_1A_2/(L^2N^2)=\zeta^2$, whereas $x(U-V)=x(2P)=v^2$.
\end{proof}
Thus nonzero guards on the two coupling coordinates do not repair
this height identity. Neither counterexample is asserted to satisfy
a complete global detector system.

\section{All quadratic identities at the fixed reference}
Let $I_{\le2}(\calV)$ be the $K$-vector space of polynomials of total
degree at most two vanishing on $\calV$ for \eqref{eq:fixedref}.
There are $66$ monomials of degree at most two in ten variables.

\begin{theorem}\label{thm:quadrics}
The space $I_{\le2}(\calV)$ has dimension $19$. Every one of its
polynomials vanishes on the whole family \eqref{eq:family}.
A basis consists of the seven nonlinear equations in
\eqref{eq:first}--\eqref{eq:coupling1}, the linear polynomial
\[
                         \lambda=e(wA_2-uA_1)-SN-TL,
\]
its ten coordinate multiples, and the reciprocal-chord polynomial
\begin{equation}\label{eq:extra}
 A_1G-\frac qr A_1F-ew\frac qr F(SN+TL).
\end{equation}
\end{theorem}
\begin{proof}
Every displayed polynomial vanishes under \eqref{eq:map}.
For \eqref{eq:extra}, use
\[
 A_1q(g-f)=ewqfLN(s+t)=ew\frac qr F(SN+TL).
\]
The other identities follow from the point equations and
Lemma~\ref{lem:line}. Every displayed polynomial also vanishes on
\eqref{eq:family}: the eight source equations do so by
Proposition~\ref{prop:family}, and \eqref{eq:extra} vanishes because
$f=g$ and $SN+TL=0$.

It remains to prove completeness of the list. Work modulo
$p=1009$ with $\rho=439$, so $\rho^2=2$. The exact certificate
of Section~\ref{sec:certificates} supplies $47$ secant pairs and a
$47$ by $47$ evaluation minor with determinant $541$ modulo $1009$.
The coefficient vectors of the nineteen displayed polynomials have
a $19$ by $19$ minor of determinant $968$ modulo $1009$.
All constants inverted in \eqref{eq:fixedref} are units at this prime,
and every sampled $fgst(g-f)$ is nonzero. In particular, the nineteen
polynomials are linearly independent in characteristic zero.

For the evaluation lower bound in characteristic zero, let
$\mathcal O=\Z[\rho]_{(1009,\rho-439)}$ and invert the fixed
reference denominators, which are units. The algebra
\[
 \mathcal O[f^{\pm1},g^{\pm1},s,t]\big/
       (s^2-f^{-1}+Df,\ t^2-g^{-1}+Dg)
\]
is free over $\mathcal O[f^{\pm1},g^{\pm1}]$ with basis
$1,s,t,st$. Its characteristic-zero generic fiber is the coordinate
algebra of an open subset of the product of two geometrically
integral curves: each square-root radicand has an odd-order pole.

Multiplying every monomial of degree at most two in
\eqref{eq:map} by
\[
                         s^4t^4(g-f)^2
\]
clears its $s,t,g-f$ denominators. Reduction by the two monic
quadratics gives a finite Laurent-coefficient matrix over
$\mathcal O$. Its rank over $K$ is the dimension of the space of
restricted degree-at-most-two functions. The evaluation matrix modulo
$p$ factors through this coefficient matrix, followed by multiplication
of its rows by nonzero denominators. Hence its characteristic-zero
rank is at least $47$. The nineteen independent identities give the
reverse bound $66-19=47$. Thus their span is all of
$I_{\le2}(\calV)$.
\end{proof}

\begin{corollary}\label{cor:noKrepair}
Adjoining any collection of quadratic $K$-polynomial identities in
these same ten coordinates, while preserving the complete geometric
secant image, cannot remove the tangent family
\eqref{eq:family}. The statement remains true if the proposed
identities originally involve reference coordinates and specialize
to polynomials at \eqref{eq:fixedref}.
\end{corollary}

\section{Quadrics in rational components}
For an affine set $Z$ over an algebraically closed field and an
integer $d\ge0$, write
\[
 C_d(Z)=\{z: f(z)=0\text{ for every }f\text{ of total degree at most }d
                         \text{ vanishing on }Z\}.
\]

\begin{lemma}\label{lem:product}
For nonempty affine sets $Z_1,Z_2$,
\[
                       C_d(Z_1\times Z_2)=C_d(Z_1)\times C_d(Z_2).
\]
\end{lemma}
\begin{proof}
Polynomials in one factor alone prove containment from left to right.
Conversely let $z_i\in C_d(Z_i)$, and let $f$ of total degree at most
$d$ vanish on $Z_1\times Z_2$. For every $y\in Z_2$, the polynomial
$f(\,\cdot\,,y)$ vanishes on $Z_1$ and still has degree at most $d$.
Thus $f(z_1,y)=0$. The polynomial $f(z_1,\,\cdot\,)$ vanishes on
$Z_2$ and has degree at most $d$, so $f(z_1,z_2)=0$.
\end{proof}

\begin{theorem}\label{thm:rational}
Write each of the ten field coordinates as $z_i=z_{i,0}+\rho z_{i,1}$.
Every rational polynomial of total degree at most two in these twenty
rational components that vanishes on the complete geometric
restriction of scalars of the secant image also vanishes at every
rational-component point of \eqref{eq:family}, for $\zeta\in K$.
This includes polynomials involving products of conjugate components.
\end{theorem}
\begin{proof}
Let $\sigma$ be the nontrivial automorphism of $K/\Q$. After extending
scalars to an algebraic closure, restriction of scalars splits into
$\calV\times\calV^\sigma$. The change from rational components to
split coordinates is the invertible linear transformation
\[
                  z_i^+=z_{i,0}+\rho z_{i,1},\qquad
                  z_i^-=z_{i,0}-\rho z_{i,1}.
\]
It preserves total degree. Theorem~\ref{thm:quadrics}, and its
conjugate, put the two components of each point of
\eqref{eq:family} in $C_2(\calV)$ and $C_2(\calV^\sigma)$.
Polynomial identity spaces here are kernels of finite coefficient
matrices, so the statement persists under field extension.
Lemma~\ref{lem:product} completes the proof.
\end{proof}

\section{Adjoining the physical rational trace}
Let $\mathcal G\subset\A_\Q^{21}$ be the closure of the complete
geometric secant image in the twenty rational components, together
with the coordinate
\[
                         \mathsf T=\Tr_{K/\Q}x(U-V).
\]
After extension of scalars, use two independent point pairs with
parameters $(f_+,s_+,g_+,t_+)$ and $(f_-,s_-,g_-,t_-)$.
The constants in the second pair are the conjugates of those in
the first. Their difference abscissas are
\begin{equation}\label{eq:splitheight}
 h_\pm=\frac{f_\pm g_\pm(s_\pm+t_\pm)^2}
                     {(g_\pm-f_\pm)^2},\qquad
                         \mathsf T=h_++h_-.
\end{equation}
These formulas follow from the two chord factors in
Lemma~\ref{lem:line}; in particular this is the physical trace,
without an additional ratio-dependent weight.

\begin{theorem}\label{thm:trace}
The space of rational polynomials of total degree at most two
vanishing on $\mathcal G$ has dimension $59$. Every such polynomial
vanishes on the rational-component tangent family
\eqref{eq:family} with \emph{arbitrary} rational $\mathsf T$.
Thus adjoining the physical trace alone does not provide a quadratic
completion of the specified embedding.
\end{theorem}
\begin{proof}
Over $K$, write the twenty-one split coordinates as
$z_1^+,\ldots,z_{10}^+,z_1^-,\ldots,z_{10}^-,\mathsf T$.
For each factor take the eight nonlinear identities in
Theorem~\ref{thm:quadrics}, its linear identity $\lambda_\pm$,
and every product of $\lambda_\pm$ with one of the twenty-one
coordinates. These sixty polynomials are true identities of the
split graph. Every one vanishes on the two false tangent families
with arbitrary $\mathsf T$, since the eight nonlinear identities
and $\lambda_\pm$ already vanish there.

There are $253$ monomials of degree at most two in twenty-one
variables. At the same prime $1009$, the trace certificate supplies
$194$ evaluations of the true split graph and a square minor of
determinant $451$. The coefficient vectors of the sixty displayed
identities have a $59$ by $59$ minor of determinant $929$.
All four point parameters are checked against their curve equations,
and all denominators in \eqref{eq:map} and \eqref{eq:splitheight}
are nonzero. Thus the displayed identities are independent in a
$59$-dimensional subspace in characteristic zero.

For the characteristic-zero evaluation bound, use the Laurent
algebra over $\mathcal O$ on the four independent curve equations
\[
 s_\pm^2=f_\pm^{-1}-D_\pm f_\pm,\qquad
 t_\pm^2=g_\pm^{-1}-D_\pm g_\pm,
 \qquad D_+=D,\quad D_-=\sigma(D).
\]
It is free of rank $16$ over the Laurent polynomial algebra in
$f_+,g_+,f_-,g_-$. A common denominator for all monomials of
degree at most two in the split coordinates and \eqref{eq:splitheight}
is
\[
 \prod_{\epsilon\in\{+,-\}}
          s_\epsilon^4t_\epsilon^4(g_\epsilon-f_\epsilon)^4.
\]
The same monic normal-form argument as in
Theorem~\ref{thm:quadrics} makes the rank at least $194$ in
characteristic zero. The $59$ independent identities make it at
most $253-59=194$. Hence the entire identity space is spanned by
the displayed polynomials and vanishes on the false tangent families
with arbitrary trace.

The change between split coordinates and rational components is
linear and invertible. The identity space descends to $\Q$ and has
the same dimension, since kernels of the finite coefficient matrices
commute with extension of scalars. A point of \eqref{eq:family} with
$\zeta\in K$ gives rational component coordinates, and every rational
choice of $\mathsf T$ is therefore left unexcluded.
\end{proof}

Theorem~\ref{thm:trace} also applies to a proposed trace multiplied
only by an input power that equals one at $\beta=1$. It does not
apply to a ratio-dependent replacement height, which would require
its own height argument. It does not cover new helper coordinates,
a different embedding, or conditions required only on a smaller
arithmetic witness family. In particular it is not a lower bound
for the number of variables in a defining quartic. These obstructions
are algebraic and use no arithmetic premise about a prime pattern,
a pole realization, the rank of $E(K)$ or the abc conjecture.

\section{Exact certificates}\label{sec:certificates}
The package
\path{negative_results/computation/tangent_completion_v48/}
contains independent implementations for both identity spaces. The first
pair of checkers uses pairs of
rational numbers for $K$ and elementary modular elimination, and
an algebraic number field, Laurent normal forms and
finite-field matrices in SymPy, respectively. The trace checkers use
the split equations and a separate SymPy reconstruction. Neither
independent checker imports producer code.

\begin{center}
\begin{tabular}{lr}
\toprule
Certificate & Value\\
\midrule
Prime and root of $2$ & $1009,\ 439$\\
Monomials of degree at most two & $66$\\
Independent secant evaluations & $47$\\
Evaluation-minor determinant modulo $1009$ & $541$\\
Independent quadratic identities & $19$\\
Identity-minor determinant modulo $1009$ & $968$\\
Rank of the affine-linear evaluation matrix & $10$\\
\midrule
Monomials after adjoining the rational trace & $253$\\
Independent true-trace evaluations & $194$\\
Trace evaluation-minor determinant & $451$\\
Independent identities with the trace & $59$\\
Trace identity-minor determinant & $929$\\
\bottomrule
\end{tabular}
\end{center}

The certificate lists the forty-seven four-tuples $(f,s,g,t)$,
the columns of the nonzero evaluation minor and the rows and columns
of the identity minor. The independent checker reconstructs all
entries, verifies the curve and nonzero conditions, and checks both
minors. It also reduces the nine generating identities to zero in
the exact characteristic-zero function field, checks the whole
false tangent family symbolically and verifies both explicit
counterexamples. The trace certificate similarly lists all $194$ eight-tuples of point
parameters and the two nonzero minors. Its independent checker verifies
the physical abscissa formula by an explicit curve-ideal identity, all
matrix entries and the false tangent family with arbitrary trace.
The product-closure argument is the proof of
Theorem~\ref{thm:rational}; no finite numerical sample substitutes for it.

\begin{thebibliography}{9}
\bibitem{Reference}
$\mathbb Z\mathbb Q$ Diophantine Project contributors,
\emph{A three-coordinate quadratic chart for an elliptic point and its translate},
\path{papers/elliptic_charts/three_coordinate_reference_20261115_v40.tex}.
\end{thebibliography}
\end{document}
